\subsection{非被約延長ファイバーからのambient Poisson二次障害}
\label{subsec:quadric-ambient-second-order-obstruction}
前節の完全等方Log例では，固定されたquadricとその上の
Poisson構造への延長ファイバーが
$\operatorname{Spec}\CC[a,b,c]/(ac,c^2)$となった。
$c$の平方零方向について，この二次障害が固定制限を外しても
残ることを示す。これは$\PP^4$上の\emph{全}Poisson変形の問題であり，
単なる相対延長イデアルの再計算ではない。

\begin{theorem}[ambientでも消えない二次障害]
\label{thm:isotropic-ambient-second-order-obstruction}
座標を$X=(u_0,u_1,z,w_0,w_1)$とし，
\[
 q=u_0^2+2u_1w_0+2zw_1,\qquad
 \xi=u_0\partial_{u_0}+u_1\partial_{u_1}-2z\partial_z,
\]
\[
 \iota_{P_0}\Omega=dq\wedge
 \iota_\xi(du_0\wedge du_1\wedge dz)
\]
とする。$E=\sum x_i\partial_{x_i}$，$G=\operatorname{Hess}(q)/2$，
\[
 C=u_1\partial_{u_1}-z\partial_z-w_0\partial_{w_0}
 +w_1\partial_{w_1},\qquad
 A=\operatorname{diag}(0,1,-1,-1,1),
\]
\[
 T=\tfrac12 AG^{-1},\qquad
 K=qT+\tfrac12 E\wedge C
\]
と置く。$[P_0]$を$\PP^4$のPoissonスキーム$\mathcal P$の点とすると，
次が成り立つ。
\begin{enumerate}
\item $K$は$\mathcal P$の$[P_0]$における非零の一次Poisson変形を与える。
これは前節の延長ファイバーのembeddedな$c$方向である。
\item 任意の斉次二次双ベクトル$R$について
\[
 P_0+\varepsilon K+\varepsilon^2R
 \quad\pmod{\varepsilon^3}
\]
が$\PP^4$上のPoisson構造になることはない。
従って$c$方向は，制限quadricや誘導Poisson構造を変えることまで
許しても二次まで延長できない。
\item 特に$\mathcal P$は$[P_0]$で滑らかではない。
全150個の斉次二次双ベクトル係数から計算した
射影Jacobi線形化の階数は$95$であり，
$T_{[P_0]}\mathcal P$の次元は$30$である。
\end{enumerate}
\end{theorem}
\begin{proof}
二次式$q$は非退化で，$A^tG+GA=0$となる。
従って$T$は反対称で，$K$は$q$をstrict Casimirとする
双ベクトルである。前節の全延長の式から，$K$は零でない
$c$方向を表し，$P_0+\varepsilon K$は
$\CC[\varepsilon]/(\varepsilon^2)$上のPoisson構造となる。

一般の斉次二次双ベクトル$B$に対して
\[
 J(B)_{ijk}=\sum_{\ell=0}^4
 \bigl(B_{i\ell}\partial_\ell B_{jk}
 +B_{j\ell}\partial_\ell B_{ki}
 +B_{k\ell}\partial_\ell B_{ij}\bigr)
\]
と定義する。この符号規約では$\PP^4$上の射影Jacobi方程式は
$\mathcal F(B)=E\wedge J(B)=0$である。
$J(P_0)=0$および$J(P_0+K)=J(K)$は直接確認できる。
従って二次補正$R$の存在は
\begin{equation}
 E\wedge\bigl(dJ_{P_0}(R)+J(K)\bigr)=0
 \label{eq:isotropic-ambient-second-obstruction}
\end{equation}
を必要とする。ここで$R$は\emph{固定quadricや固定制限の条件を
一切課していない}任意の150係数の双ベクトルである。

\noindent
次数四の四重ベクトル$F$の成分を$F_{ijkl}$と書き，
$[m]F_{ijkl}$を単項式$m$の係数とする。次の五係数汎関数を取る。
\[
\begin{split}
 \Lambda(F)={}&2[u_0^3z]F_{0123}
 +[u_0u_1zw_0]F_{0123}+[u_0z^2w_1]F_{0123}\\
 &+4[u_0^3u_1]F_{0124}
 +10[u_0u_1^2w_0]F_{0124}.
\end{split}
\]
任意の$R$に対し
\begin{equation}
 \Lambda\bigl(E\wedge dJ_{P_0}(R)\bigr)=0
 \label{eq:isotropic-five-term-cancellation}
\end{equation}
となる。実際，150個の基底
$x_ax_b\partial_i\wedge\partial_j$上で計算すると，
この五係数に現れる非零列は六列だけで，係数行列は
\[
\begin{pmatrix}
 2&0&-12&0&2&0\\
 -2&12&0&-12&4&-6\\
 -2&0&0&0&-8&6\\
 0&2&6&-2&0&0\\
 0&-2&0&2&0&0
\end{pmatrix}.
\]
その左から$(2,1,1,4,10)$を掛けると零になる。
これは有限標本ではなく，\emph{全150係数}に対する係数恒等式である。
一方，$J(K)$の非零成分は
\[
\begin{aligned}
 J(K)_{123}&=\tfrac12zq,&
 J(K)_{124}&=-\tfrac12u_1q,\\
 J(K)_{134}&=\tfrac12w_1q,&
 J(K)_{234}&=-\tfrac12w_0q.
\end{aligned}
\]
従って，上の五係数に対する$E\wedge J(K)$の値は順に
$(\frac12,1,1,-\frac12,-1)$であり，
\[
 \Lambda\bigl(E\wedge J(K)\bigr)
 =2\cdot\tfrac12+1+1+4\cdot(-\tfrac12)+10\cdot(-1)=-9\ne0.
\]
式\eqref{eq:isotropic-ambient-second-obstruction}は
\eqref{eq:isotropic-five-term-cancellation}と矛盾する。
これで第2項が示された。

全350個の射影Jacobi係数を150個の斉次二次係数で線形化すると
階数$95$となる。$E\wedge$一次ベクトル場の24次元のゲージと
射影尺度の一次元を引けば，接空間次元は$150-95-24-1=30$である。
この係数階数と五係数恒等式は，付属する独立の
\texttt{check\_ambient\_quadratic\_obstruction.py}で$\mathbb Q$上に再構成できる。
滑らかなスキームでは任意の一次変形が
$\CC[\varepsilon]/(\varepsilon^3)$まで持ち上がるので，
上の非持上げ性は$[P_0]$の非滑らかさを含意する。
\end{proof}

\begin{remark}[非被約性と特異性の区別]
\label{rem:ambient-obstruction-not-nonreduced}
本定理はambient Poissonスキームの\emph{二次障害が非零}であることを
証明するが，ambientの局所環が非被約であることは証明しない。
被約な特異スキームにも二次まで持ち上がらない接方向はある。
固定quadricの延長ファイバーの非被約性は前節の別の定理であり，
両者を混同しない。なお，被約台上の一般rank四点が
既存147成分のどれに属するかは，ここでは決定していない。
\end{remark}


\begin{corollary}[完全等方Log点での二次接錐と局所次元]
\label{cor:ambient-isotropic-tangent-cone-bound}
定理\ref{thm:isotropic-ambient-second-order-obstruction}の
$p=[P_0]\in\mathcal P=\mathbb P\operatorname{Pois}(\PP^4)$に対し、
既述の五係数汎関数$\Lambda$から定まる二次式
\[
 \mathcal Q_\Lambda(B)=\Lambda\bigl(E\wedge J(B)\bigr)
\]
は$30$次元のZariski接空間$T_p\mathcal P$上の非零二次形式を定め、
$p$での接錐はその零集合に含まれる。従って
\[
 \boxed{\dim T_p\mathcal P=30,\qquad\dim_p\mathcal P\le29.}
\]
特に$p$を含む既約成分の次元はすべて$29$以下である。
現行カタログの次元$36$の\texttt{R2\_RAT113}と
次元$39$の\texttt{R2\_LPB}は$p$を含まない。

さらに完全等方Logの固定$Q$延長ファイバーを表す場
\[
 A_1=u_1\partial_{u_0}-u_0\partial_{w_0},\quad
 A_2=z\partial_{w_0}-u_1\partial_{w_1},\quad A_3=C
\]
に対応する式\eqref{eq:isotropic-K}の核双ベクトルを
$K_1,K_2,K_3$とすると、
\begin{equation}\label{eq:isotropic-tangent-slice}
 \mathcal Q_\Lambda(P_0+aK_1+bK_2+cK_3)=-9c^2
\end{equation}
が$a,b,c$の恒等式として成立する。
\end{corollary}
\begin{proof}
全$150$係数で成立する五係数消滅恒等式
\eqref{eq:isotropic-five-term-cancellation}は
$\Lambda(d\mathcal F_{P_0}(B))=0$を意味する。
$\mathcal F$は二次で$\mathcal F(P_0)=0$なので、
局所座標$P_0+B$では
$\Lambda(\mathcal F(P_0+B))=\mathcal Q_\Lambda(B)$となる。

任意の線形ベクトル場$V$に対し、
$[E,B]=[E,V]=0$とSchouten括弧のLeibniz則から
$E\wedge J(B+E\wedge V)=E\wedge J(B)$である。
また上述の線形化の消滅から
$\mathcal Q_\Lambda(B+\lambda P_0)=\mathcal Q_\Lambda(B)$となる。
従ってこの二次式はEulerゲージおよび射影尺度の商である
$T_p\mathcal P$に降りる。
前定理の接方向$K=K_3$で
$\mathcal Q_\Lambda(K)=-9\ne0$なので、
接錐の定義イデアルには接空間上の非零二次式が含まれる。
局所環と接錐のKrull次元は同じであるから
$\dim_p\mathcal P\le30-1=29$を得る。
カタログの二成分は次元がこの上界を超えるので$p$を含まない。

最後に$A_1,A_2,A_3$を式\eqref{eq:isotropic-K}に代入し、
$E\wedge J(P_0+aK_1+bK_2+cK_3)$の
$\Lambda$で抽出する五係数を直接展開すると
\[
 (c^2/2,\ c^2,\ c^2,\ -c^2/2,\ -c^2)
\]
となる。これに重み$(2,1,1,4,10)$を掛ければ$-9c^2$。
この恒等式は元の障害producerとは独立にSymPyで再計算した。
\end{proof}
\begin{remark}[未解決の局所幾何]
この系は局所次元の等号成立、接錐の既約分解、
$p$を含む残る成分の特定、ambient局所環の非被約性を意味しない。
表示した二次式が検出するのは固定ファイバーの$c^2$関係であり、
$ac$関係のambient障害性も未確認である。
\end{remark}
